
\chapter*{Abstract}
Automated abnormality detection from 3D chest Computed Tomography requires clinically interpretable methods to support radiological
decision-making.
Current State-of-the-Art (SotA) deep learning models often rely on post-hoc explainability techniques that provide approximate localizations
and lack direct traceability to the models decisions.
Furthermore, SotA models rarely support simultaneous multi-class classification and grounded localization, leading to a fragmented
approach to diagnosis. This also hinders the ability to benchmark models across multiple dimensions of performance,
including classification accuracy, segmentation quality, and explainability.
This study aims to address these challenges through two contributions:
First, we introduce an inherently interpretable dictionary learning approach in which detected abnormalities can be
back-projected to the original voxel space, allowing diagnostic decisions to be traced back to relevant anatomical structures.
Second, we develop a unified, open-source evaluation framework to benchmark SotA models across three axes: (i) classification accuracy,
(ii) segmentation quality, and (iii) explainability, using standardized metrics and a publically available 3D chest CT dataset.
Our method is benchmarked against SotA models using this framework, demonstrating competitive
classification and segmentation performance while providing improved interpretability compared to deep learning baselines.
